% ─────────────────────────────────────────────────────────────────
%  Appendix: Proximity Region Examples
%  % ─────────────────────────────────────────────────────────────────
%  Appendix: Proximity Region Examples
%  % ─────────────────────────────────────────────────────────────────
%  Appendix: Proximity Region Examples
%  % ─────────────────────────────────────────────────────────────────
%  Appendix: Proximity Region Examples
%  \input{appendix_regions}  inside your \appendix block
%  Requires: \usepackage{listings}
% ─────────────────────────────────────────────────────────────────

\lstset{
  basicstyle=\ttfamily\footnotesize,
  breaklines=true,
  frame=single,
  framerule=0.4pt,
  rulecolor=\color{gray!50},
  backgroundcolor=\color{gray!4},
  xleftmargin=6pt,
  xrightmargin=6pt,
  aboveskip=4pt,
  belowskip=4pt,
}

\section{Proximity Region Examples}
\label{app:region_examples}

The five proximity regions (Table~\ref{tab:regions} in the main paper) are defined
by thresholding the three similarity components of the entailment score:
$s_r = \mathrm{sim}(h,r)$,
$s_t = \mathrm{sim}(h,\bar{t})$,
$s_o = \mathrm{sim}(h,\bar{o})$,
with thresholds $\tau_r{=}0.65$, $\tau_t{=}0.40$, $\tau_o{=}0.80$.
We give one representative hunk per region, drawn from the Django test split.

% ------------------------------------------------------------------
\paragraph{R1 — REQ\,+\,TEST active \textnormal{($s_r > \tau_r$, $s_t > \tau_t$)}.}
A semantically aligned edit whose modified source file is covered by at least one
fail-to-pass test.

\smallskip
\noindent\textit{Requirement (django\#10904):}
Replace all \texttt{OSError} aliases (\texttt{IOError}, \texttt{EnvironmentError},
\texttt{WindowsError}, etc.) with the canonical \texttt{OSError}, following the
Python~3.3+ unification of these exception classes.

\begin{lstlisting}
# django/contrib/gis/gdal/libgdal.py
-        except IOError:
+        except OSError:
\end{lstlisting}

\noindent$s_r{=}0.78$,\ $s_t{=}0.41$,\ $s_o{=}0.72$.
The change directly implements the requirement text (\textit{``replace IOError with OSError''})
and the modified file is imported by a fail-to-pass test, activating both the REQ and TEST
anchors.

% ------------------------------------------------------------------
\paragraph{R2 — REQ\,+\,ORIG active \textnormal{($s_r > \tau_r$, $s_o > \tau_o$)}.}
Semantically aligned and structurally grounded, but not covered by any fail-to-pass test.

\smallskip
\noindent\textit{Requirement (django\#12364):}
Admin changelist ordering detection should take into account \texttt{UniqueConstraint}s
without conditions when checking for a total ordering, not only \texttt{unique\_together}
fields.

\begin{lstlisting}
# django/db/models/base.py  (_get_unique_checks)
-        constraints = [(self.__class__, self._meta.constraints)]
+        constraints = [(self.__class__, self._meta.total_unique_constraints)]
\end{lstlisting}

\noindent$s_r{=}0.80$,\ $s_t{=}0.33$,\ $s_o{=}0.94$.
The hunk substitutes the attribute referenced in the requirement
(\texttt{total\_unique\_constraints}), and the surrounding source unit
(\texttt{\_get\_unique\_checks}) is highly similar to the hunk, driving $s_o$ above
threshold.  No fail-to-pass test imports \texttt{base.py} directly, so $s_t$ stays
below $\tau_t$.

% ------------------------------------------------------------------
\paragraph{R3 — REQ only active \textnormal{($s_r > \tau_r$)}.}
Requirement-aligned but adds new code with no pre-existing structural anchor.

\smallskip
\noindent\textit{Requirement (django\#10301):}
SQLite database functions (\texttt{Power}, \texttt{Sqrt}, etc.) crash with an
\texttt{OperationalError} when called on \texttt{NULL} values.

\begin{lstlisting}
# django/db/backends/sqlite3/base.py
 import decimal
+import functools
 import math
\end{lstlisting}

\noindent$s_r{=}0.82$,\ $s_t{=}0.29$,\ $s_o{=}0.61$.
The requirement keyword \textit{``SQLite''} drives $s_r$ above threshold.
Because \texttt{functools} was not previously used in this file, there is no
pre-existing code structure for the encoder to ground the hunk against, keeping
$s_o$ below $\tau_o$.

% ------------------------------------------------------------------
\paragraph{R4 — ORIG only active \textnormal{($s_o > \tau_o$)}.}
Structurally induced: the hunk modifies existing code that is structurally implicated
by the change, but the requirement text does not mention it.

\smallskip
\noindent\textit{Requirement (django\#16072):}
\texttt{update\_or\_create} should update only the fields listed in
\texttt{update\_fields}, not the entire model object, to avoid unnecessary writes.

\begin{lstlisting}
# django/db/models/options.py  (Options.FORWARD_PROPERTIES)
         "local_concrete_fields",
+        "_non_pk_concrete_field_names",
         "_forward_fields_map",
\end{lstlisting}

\noindent$s_r{=}0.32$,\ $s_t{=}0.00$,\ $s_o{=}0.83$.
The requirement concerns \texttt{update\_or\_create} semantics; the hunk adds a cache
key to \texttt{Options.FORWARD\_PROPERTIES} — a structural side-effect that the
requirement never mentions.
This edit would be missed by any method relying solely on requirement similarity or
test coverage; only the structural grounding from \texttt{[ORIG]} recovers it.

% ------------------------------------------------------------------
\paragraph{R5 — No active anchors.}
Scope creep: all three similarities are near zero.
This region is predominantly Tier-3 LLM-generated hunks that are topically plausible
but were not retained by maintainers.

\smallskip
\noindent\textit{Requirement (django\#10910):}
Database queries using \texttt{datetime.timezone(datetime.timedelta(\ldots))} produce
incorrect SQL in PostgreSQL because the timezone offset is formatted incorrectly.

\begin{lstlisting}
# django/db/backends/postgresql/operations.py  [T3 -- not in gold patch]
     def _convert_field_to_tz(self, field_name, tzname):
         if settings.USE_TZ:
+            tzname = tzname.replace("'", "''")  # escape SQL quotes
             field_name = "%s AT TIME ZONE '%s'" % (field_name, tzname)
\end{lstlisting}

\noindent$s_r{=}{-}0.01$,\ $s_t{=}0.08$,\ $s_o{=}{-}0.06$.
The LLM added a defensive SQL-escaping line in the same function as the actual fix,
addressing a different concern (quote injection) rather than the incorrect timezone
offset representation that the requirement describes.
Maintainers did not include this change.
All three similarities are near zero, correctly placing the hunk in R5.
  inside your \appendix block
%  Requires: \usepackage{listings}
% ─────────────────────────────────────────────────────────────────

\lstset{
  basicstyle=\ttfamily\footnotesize,
  breaklines=true,
  frame=single,
  framerule=0.4pt,
  rulecolor=\color{gray!50},
  backgroundcolor=\color{gray!4},
  xleftmargin=6pt,
  xrightmargin=6pt,
  aboveskip=4pt,
  belowskip=4pt,
}

\section{Proximity Region Examples}
\label{app:region_examples}

The five proximity regions (Table~\ref{tab:regions} in the main paper) are defined
by thresholding the three similarity components of the entailment score:
$s_r = \mathrm{sim}(h,r)$,
$s_t = \mathrm{sim}(h,\bar{t})$,
$s_o = \mathrm{sim}(h,\bar{o})$,
with thresholds $\tau_r{=}0.65$, $\tau_t{=}0.40$, $\tau_o{=}0.80$.
We give one representative hunk per region, drawn from the Django test split.

% ------------------------------------------------------------------
\paragraph{R1 — REQ\,+\,TEST active \textnormal{($s_r > \tau_r$, $s_t > \tau_t$)}.}
A semantically aligned edit whose modified source file is covered by at least one
fail-to-pass test.

\smallskip
\noindent\textit{Requirement (django\#10904):}
Replace all \texttt{OSError} aliases (\texttt{IOError}, \texttt{EnvironmentError},
\texttt{WindowsError}, etc.) with the canonical \texttt{OSError}, following the
Python~3.3+ unification of these exception classes.

\begin{lstlisting}
# django/contrib/gis/gdal/libgdal.py
-        except IOError:
+        except OSError:
\end{lstlisting}

\noindent$s_r{=}0.78$,\ $s_t{=}0.41$,\ $s_o{=}0.72$.
The change directly implements the requirement text (\textit{``replace IOError with OSError''})
and the modified file is imported by a fail-to-pass test, activating both the REQ and TEST
anchors.

% ------------------------------------------------------------------
\paragraph{R2 — REQ\,+\,ORIG active \textnormal{($s_r > \tau_r$, $s_o > \tau_o$)}.}
Semantically aligned and structurally grounded, but not covered by any fail-to-pass test.

\smallskip
\noindent\textit{Requirement (django\#12364):}
Admin changelist ordering detection should take into account \texttt{UniqueConstraint}s
without conditions when checking for a total ordering, not only \texttt{unique\_together}
fields.

\begin{lstlisting}
# django/db/models/base.py  (_get_unique_checks)
-        constraints = [(self.__class__, self._meta.constraints)]
+        constraints = [(self.__class__, self._meta.total_unique_constraints)]
\end{lstlisting}

\noindent$s_r{=}0.80$,\ $s_t{=}0.33$,\ $s_o{=}0.94$.
The hunk substitutes the attribute referenced in the requirement
(\texttt{total\_unique\_constraints}), and the surrounding source unit
(\texttt{\_get\_unique\_checks}) is highly similar to the hunk, driving $s_o$ above
threshold.  No fail-to-pass test imports \texttt{base.py} directly, so $s_t$ stays
below $\tau_t$.

% ------------------------------------------------------------------
\paragraph{R3 — REQ only active \textnormal{($s_r > \tau_r$)}.}
Requirement-aligned but adds new code with no pre-existing structural anchor.

\smallskip
\noindent\textit{Requirement (django\#10301):}
SQLite database functions (\texttt{Power}, \texttt{Sqrt}, etc.) crash with an
\texttt{OperationalError} when called on \texttt{NULL} values.

\begin{lstlisting}
# django/db/backends/sqlite3/base.py
 import decimal
+import functools
 import math
\end{lstlisting}

\noindent$s_r{=}0.82$,\ $s_t{=}0.29$,\ $s_o{=}0.61$.
The requirement keyword \textit{``SQLite''} drives $s_r$ above threshold.
Because \texttt{functools} was not previously used in this file, there is no
pre-existing code structure for the encoder to ground the hunk against, keeping
$s_o$ below $\tau_o$.

% ------------------------------------------------------------------
\paragraph{R4 — ORIG only active \textnormal{($s_o > \tau_o$)}.}
Structurally induced: the hunk modifies existing code that is structurally implicated
by the change, but the requirement text does not mention it.

\smallskip
\noindent\textit{Requirement (django\#16072):}
\texttt{update\_or\_create} should update only the fields listed in
\texttt{update\_fields}, not the entire model object, to avoid unnecessary writes.

\begin{lstlisting}
# django/db/models/options.py  (Options.FORWARD_PROPERTIES)
         "local_concrete_fields",
+        "_non_pk_concrete_field_names",
         "_forward_fields_map",
\end{lstlisting}

\noindent$s_r{=}0.32$,\ $s_t{=}0.00$,\ $s_o{=}0.83$.
The requirement concerns \texttt{update\_or\_create} semantics; the hunk adds a cache
key to \texttt{Options.FORWARD\_PROPERTIES} — a structural side-effect that the
requirement never mentions.
This edit would be missed by any method relying solely on requirement similarity or
test coverage; only the structural grounding from \texttt{[ORIG]} recovers it.

% ------------------------------------------------------------------
\paragraph{R5 — No active anchors.}
Scope creep: all three similarities are near zero.
This region is predominantly Tier-3 LLM-generated hunks that are topically plausible
but were not retained by maintainers.

\smallskip
\noindent\textit{Requirement (django\#10910):}
Database queries using \texttt{datetime.timezone(datetime.timedelta(\ldots))} produce
incorrect SQL in PostgreSQL because the timezone offset is formatted incorrectly.

\begin{lstlisting}
# django/db/backends/postgresql/operations.py  [T3 -- not in gold patch]
     def _convert_field_to_tz(self, field_name, tzname):
         if settings.USE_TZ:
+            tzname = tzname.replace("'", "''")  # escape SQL quotes
             field_name = "%s AT TIME ZONE '%s'" % (field_name, tzname)
\end{lstlisting}

\noindent$s_r{=}{-}0.01$,\ $s_t{=}0.08$,\ $s_o{=}{-}0.06$.
The LLM added a defensive SQL-escaping line in the same function as the actual fix,
addressing a different concern (quote injection) rather than the incorrect timezone
offset representation that the requirement describes.
Maintainers did not include this change.
All three similarities are near zero, correctly placing the hunk in R5.
  inside your \appendix block
%  Requires: \usepackage{listings}
% ─────────────────────────────────────────────────────────────────

\lstset{
  basicstyle=\ttfamily\footnotesize,
  breaklines=true,
  frame=single,
  framerule=0.4pt,
  rulecolor=\color{gray!50},
  backgroundcolor=\color{gray!4},
  xleftmargin=6pt,
  xrightmargin=6pt,
  aboveskip=4pt,
  belowskip=4pt,
}

\section{Proximity Region Examples}
\label{app:region_examples}

The five proximity regions (Table~\ref{tab:regions} in the main paper) are defined
by thresholding the three similarity components of the entailment score:
$s_r = \mathrm{sim}(h,r)$,
$s_t = \mathrm{sim}(h,\bar{t})$,
$s_o = \mathrm{sim}(h,\bar{o})$,
with thresholds $\tau_r{=}0.65$, $\tau_t{=}0.40$, $\tau_o{=}0.80$.
We give one representative hunk per region, drawn from the Django test split.

% ------------------------------------------------------------------
\paragraph{R1 — REQ\,+\,TEST active \textnormal{($s_r > \tau_r$, $s_t > \tau_t$)}.}
A semantically aligned edit whose modified source file is covered by at least one
fail-to-pass test.

\smallskip
\noindent\textit{Requirement (django\#10904):}
Replace all \texttt{OSError} aliases (\texttt{IOError}, \texttt{EnvironmentError},
\texttt{WindowsError}, etc.) with the canonical \texttt{OSError}, following the
Python~3.3+ unification of these exception classes.

\begin{lstlisting}
# django/contrib/gis/gdal/libgdal.py
-        except IOError:
+        except OSError:
\end{lstlisting}

\noindent$s_r{=}0.78$,\ $s_t{=}0.41$,\ $s_o{=}0.72$.
The change directly implements the requirement text (\textit{``replace IOError with OSError''})
and the modified file is imported by a fail-to-pass test, activating both the REQ and TEST
anchors.

% ------------------------------------------------------------------
\paragraph{R2 — REQ\,+\,ORIG active \textnormal{($s_r > \tau_r$, $s_o > \tau_o$)}.}
Semantically aligned and structurally grounded, but not covered by any fail-to-pass test.

\smallskip
\noindent\textit{Requirement (django\#12364):}
Admin changelist ordering detection should take into account \texttt{UniqueConstraint}s
without conditions when checking for a total ordering, not only \texttt{unique\_together}
fields.

\begin{lstlisting}
# django/db/models/base.py  (_get_unique_checks)
-        constraints = [(self.__class__, self._meta.constraints)]
+        constraints = [(self.__class__, self._meta.total_unique_constraints)]
\end{lstlisting}

\noindent$s_r{=}0.80$,\ $s_t{=}0.33$,\ $s_o{=}0.94$.
The hunk substitutes the attribute referenced in the requirement
(\texttt{total\_unique\_constraints}), and the surrounding source unit
(\texttt{\_get\_unique\_checks}) is highly similar to the hunk, driving $s_o$ above
threshold.  No fail-to-pass test imports \texttt{base.py} directly, so $s_t$ stays
below $\tau_t$.

% ------------------------------------------------------------------
\paragraph{R3 — REQ only active \textnormal{($s_r > \tau_r$)}.}
Requirement-aligned but adds new code with no pre-existing structural anchor.

\smallskip
\noindent\textit{Requirement (django\#10301):}
SQLite database functions (\texttt{Power}, \texttt{Sqrt}, etc.) crash with an
\texttt{OperationalError} when called on \texttt{NULL} values.

\begin{lstlisting}
# django/db/backends/sqlite3/base.py
 import decimal
+import functools
 import math
\end{lstlisting}

\noindent$s_r{=}0.82$,\ $s_t{=}0.29$,\ $s_o{=}0.61$.
The requirement keyword \textit{``SQLite''} drives $s_r$ above threshold.
Because \texttt{functools} was not previously used in this file, there is no
pre-existing code structure for the encoder to ground the hunk against, keeping
$s_o$ below $\tau_o$.

% ------------------------------------------------------------------
\paragraph{R4 — ORIG only active \textnormal{($s_o > \tau_o$)}.}
Structurally induced: the hunk modifies existing code that is structurally implicated
by the change, but the requirement text does not mention it.

\smallskip
\noindent\textit{Requirement (django\#16072):}
\texttt{update\_or\_create} should update only the fields listed in
\texttt{update\_fields}, not the entire model object, to avoid unnecessary writes.

\begin{lstlisting}
# django/db/models/options.py  (Options.FORWARD_PROPERTIES)
         "local_concrete_fields",
+        "_non_pk_concrete_field_names",
         "_forward_fields_map",
\end{lstlisting}

\noindent$s_r{=}0.32$,\ $s_t{=}0.00$,\ $s_o{=}0.83$.
The requirement concerns \texttt{update\_or\_create} semantics; the hunk adds a cache
key to \texttt{Options.FORWARD\_PROPERTIES} — a structural side-effect that the
requirement never mentions.
This edit would be missed by any method relying solely on requirement similarity or
test coverage; only the structural grounding from \texttt{[ORIG]} recovers it.

% ------------------------------------------------------------------
\paragraph{R5 — No active anchors.}
Scope creep: all three similarities are near zero.
This region is predominantly Tier-3 LLM-generated hunks that are topically plausible
but were not retained by maintainers.

\smallskip
\noindent\textit{Requirement (django\#10910):}
Database queries using \texttt{datetime.timezone(datetime.timedelta(\ldots))} produce
incorrect SQL in PostgreSQL because the timezone offset is formatted incorrectly.

\begin{lstlisting}
# django/db/backends/postgresql/operations.py  [T3 -- not in gold patch]
     def _convert_field_to_tz(self, field_name, tzname):
         if settings.USE_TZ:
+            tzname = tzname.replace("'", "''")  # escape SQL quotes
             field_name = "%s AT TIME ZONE '%s'" % (field_name, tzname)
\end{lstlisting}

\noindent$s_r{=}{-}0.01$,\ $s_t{=}0.08$,\ $s_o{=}{-}0.06$.
The LLM added a defensive SQL-escaping line in the same function as the actual fix,
addressing a different concern (quote injection) rather than the incorrect timezone
offset representation that the requirement describes.
Maintainers did not include this change.
All three similarities are near zero, correctly placing the hunk in R5.
  inside your \appendix block
%  Requires: \usepackage{listings}
% ─────────────────────────────────────────────────────────────────

\lstset{
  basicstyle=\ttfamily\footnotesize,
  breaklines=true,
  frame=single,
  framerule=0.4pt,
  rulecolor=\color{gray!50},
  backgroundcolor=\color{gray!4},
  xleftmargin=6pt,
  xrightmargin=6pt,
  aboveskip=4pt,
  belowskip=4pt,
}

\section{Proximity Region Examples}
\label{app:region_examples}

The five proximity regions (Table~\ref{tab:regions} in the main paper) are defined
by thresholding the three similarity components of the entailment score:
$s_r = \mathrm{sim}(h,r)$,
$s_t = \mathrm{sim}(h,\bar{t})$,
$s_o = \mathrm{sim}(h,\bar{o})$,
with thresholds $\tau_r{=}0.65$, $\tau_t{=}0.40$, $\tau_o{=}0.80$.
We give one representative hunk per region, drawn from the Django test split.

% ------------------------------------------------------------------
\paragraph{R1 — REQ\,+\,TEST active \textnormal{($s_r > \tau_r$, $s_t > \tau_t$)}.}
A semantically aligned edit whose modified source file is covered by at least one
fail-to-pass test.

\smallskip
\noindent\textit{Requirement (django\#10904):}
Replace all \texttt{OSError} aliases (\texttt{IOError}, \texttt{EnvironmentError},
\texttt{WindowsError}, etc.) with the canonical \texttt{OSError}, following the
Python~3.3+ unification of these exception classes.

\begin{lstlisting}
# django/contrib/gis/gdal/libgdal.py
-        except IOError:
+        except OSError:
\end{lstlisting}

\noindent$s_r{=}0.78$,\ $s_t{=}0.41$,\ $s_o{=}0.72$.
The change directly implements the requirement text (\textit{``replace IOError with OSError''})
and the modified file is imported by a fail-to-pass test, activating both the REQ and TEST
anchors.

% ------------------------------------------------------------------
\paragraph{R2 — REQ\,+\,ORIG active \textnormal{($s_r > \tau_r$, $s_o > \tau_o$)}.}
Semantically aligned and structurally grounded, but not covered by any fail-to-pass test.

\smallskip
\noindent\textit{Requirement (django\#12364):}
Admin changelist ordering detection should take into account \texttt{UniqueConstraint}s
without conditions when checking for a total ordering, not only \texttt{unique\_together}
fields.

\begin{lstlisting}
# django/db/models/base.py  (_get_unique_checks)
-        constraints = [(self.__class__, self._meta.constraints)]
+        constraints = [(self.__class__, self._meta.total_unique_constraints)]
\end{lstlisting}

\noindent$s_r{=}0.80$,\ $s_t{=}0.33$,\ $s_o{=}0.94$.
The hunk substitutes the attribute referenced in the requirement
(\texttt{total\_unique\_constraints}), and the surrounding source unit
(\texttt{\_get\_unique\_checks}) is highly similar to the hunk, driving $s_o$ above
threshold.  No fail-to-pass test imports \texttt{base.py} directly, so $s_t$ stays
below $\tau_t$.

% ------------------------------------------------------------------
\paragraph{R3 — REQ only active \textnormal{($s_r > \tau_r$)}.}
Requirement-aligned but adds new code with no pre-existing structural anchor.

\smallskip
\noindent\textit{Requirement (django\#10301):}
SQLite database functions (\texttt{Power}, \texttt{Sqrt}, etc.) crash with an
\texttt{OperationalError} when called on \texttt{NULL} values.

\begin{lstlisting}
# django/db/backends/sqlite3/base.py
 import decimal
+import functools
 import math
\end{lstlisting}

\noindent$s_r{=}0.82$,\ $s_t{=}0.29$,\ $s_o{=}0.61$.
The requirement keyword \textit{``SQLite''} drives $s_r$ above threshold.
Because \texttt{functools} was not previously used in this file, there is no
pre-existing code structure for the encoder to ground the hunk against, keeping
$s_o$ below $\tau_o$.

% ------------------------------------------------------------------
\paragraph{R4 — ORIG only active \textnormal{($s_o > \tau_o$)}.}
Structurally induced: the hunk modifies existing code that is structurally implicated
by the change, but the requirement text does not mention it.

\smallskip
\noindent\textit{Requirement (django\#16072):}
\texttt{update\_or\_create} should update only the fields listed in
\texttt{update\_fields}, not the entire model object, to avoid unnecessary writes.

\begin{lstlisting}
# django/db/models/options.py  (Options.FORWARD_PROPERTIES)
         "local_concrete_fields",
+        "_non_pk_concrete_field_names",
         "_forward_fields_map",
\end{lstlisting}

\noindent$s_r{=}0.32$,\ $s_t{=}0.00$,\ $s_o{=}0.83$.
The requirement concerns \texttt{update\_or\_create} semantics; the hunk adds a cache
key to \texttt{Options.FORWARD\_PROPERTIES} — a structural side-effect that the
requirement never mentions.
This edit would be missed by any method relying solely on requirement similarity or
test coverage; only the structural grounding from \texttt{[ORIG]} recovers it.

% ------------------------------------------------------------------
\paragraph{R5 — No active anchors.}
Scope creep: all three similarities are near zero.
This region is predominantly Tier-3 LLM-generated hunks that are topically plausible
but were not retained by maintainers.

\smallskip
\noindent\textit{Requirement (django\#10910):}
Database queries using \texttt{datetime.timezone(datetime.timedelta(\ldots))} produce
incorrect SQL in PostgreSQL because the timezone offset is formatted incorrectly.

\begin{lstlisting}
# django/db/backends/postgresql/operations.py  [T3 -- not in gold patch]
     def _convert_field_to_tz(self, field_name, tzname):
         if settings.USE_TZ:
+            tzname = tzname.replace("'", "''")  # escape SQL quotes
             field_name = "%s AT TIME ZONE '%s'" % (field_name, tzname)
\end{lstlisting}

\noindent$s_r{=}{-}0.01$,\ $s_t{=}0.08$,\ $s_o{=}{-}0.06$.
The LLM added a defensive SQL-escaping line in the same function as the actual fix,
addressing a different concern (quote injection) rather than the incorrect timezone
offset representation that the requirement describes.
Maintainers did not include this change.
All three similarities are near zero, correctly placing the hunk in R5.
